\section{Experimental Setup}
\label{sec:experiments}

\paragraph{Models.}
We evaluate three runs:
\begin{itemize}[noitemsep]
  \item \textbf{E1 — Qwen2.5-Math-7B-Instruct}~\citep{qwen2025math}:
    a math-specialized 7B model loaded in bf16 on a single NVIDIA A100-SXM4
    (40\,GB).
    Used for all four REFUTE strategies, conjecture refinement, and
    Strategist LLM fallback.
  \item \textbf{E2 — Qwen2.5-32B-Instruct}~\citep{qwen2025}:
    a general-purpose 32B model loaded in 4-bit NF4 quantization via
    bitsandbytes~0.49~\citep{dettmers2023qlora} on the same A100-40\,GB
    (${\approx}19$\,GB VRAM at inference).
    Identical settings to E1 except model and concurrency.
  \item \textbf{E3 — Qwen2.5-Math-7B-Instruct (patched verifier)}:
    identical to E1 in all settings, but run after the SymPy NL-parser fix
    described in \S\ref{sec:analysis}.
    E3 isolates the recall impact of the verifier bug.
\end{itemize}

\paragraph{Inference.}
Both models are served via \texttt{transformers} \texttt{pipeline} with
\texttt{device\_map=auto} and \texttt{max\_new\_tokens=4096}.
No model is fine-tuned; all results reflect zero-shot or few-shot prompting
as described in \S\ref{sec:system}.

\paragraph{REFUTE configuration.}
Each conjecture is allocated \texttt{max\_rounds=10} R-Agent rounds and up to
\texttt{max\_refinements=3} C-Agent refinements.
E1 uses \texttt{max\_concurrent=2}; E2 uses \texttt{max\_concurrent=1} due to
the larger model footprint.
Both runs use \texttt{seed=0} for Python \texttt{random} seeding
(LLM outputs and structured-random candidates remain non-deterministic).

\paragraph{Evaluation metrics.}
Precision, recall, and F$_1$ are computed over the binary refutation
decision (refuted vs.\ survived) against ground-truth falsity labels.
A conjecture with status \texttt{refined} is treated as refuted
(the original statement was disproved).
Bootstrap 95\% confidence intervals use $n=1{,}000$ resamples.
